% !Mode:: "TeX:UTF-8"

\chapter{数据集构造与实验设计}

\section{空间场景生成}

实验以楼梯转角和电动轮椅几何参数为核心，生成多个空间场景。每个场景包含楼梯净宽、平台进深、轮椅宽度、轮椅长度、最小转弯半径等字段。项目中使用的参考余量定义为：

\begin{equation}
margin = \min(W - w,\ D - R)
\label{eq:margin}
\end{equation}

其中，$W$ 表示楼梯通道净宽，$w$ 表示轮椅宽度，$D$ 表示平台有效进深，$R$ 表示保守转弯半径。该参考余量不直接告诉模型，而是用于实验标注与后续评价。

\section{Prompt 描述层级}

实验将同一空间场景转写为三个描述层级，如表 \ref{tab:levels} 所示。该设计用于比较模型在自然语言直觉、参数比较和形式化几何建模三种语境下的行为差异。

\begin{table}[htbp]
    \centering
    \bicaption{Prompt 描述层级设计}{Design of prompt description levels}
    \begin{tabular}{lll}
        \toprule
        层级 & 描述方式 & 实验目的 \\
        \midrule
        L1 & 自然语言描述 & 测试模型是否依赖空间常识和直觉判断 \\
        L2 & 具体参数描述 & 测试模型能否正确使用尺寸和阈值信息 \\
        L3 & 数学建模描述 & 测试模型面对形式化几何约束时的稳定性 \\
        \bottomrule
    \end{tabular}
    \label{tab:levels}
\end{table}

\section{干扰条件设计}

为了检验模型是否会被无关环境信息带偏，实验为每个 Prompt 设置 5 类干扰条件。干扰项本身不改变空间几何约束，但可能影响模型的语义注意力分配。

\begin{table}[htbp]
    \centering
    \bicaption{无关干扰条件设计}{Design of irrelevant perturbation conditions}
    \begin{tabular}{ll}
        \toprule
        干扰标签 & 示例内容 \\
        \midrule
        none & 不添加无关信息 \\
        material\_color & 楼梯扶手颜色、墙面材质 \\
        ambient\_noise & 现场说话声和环境噪声 \\
        lighting & 光照偏暗、地面干燥等状态 \\
        emotional\_pressure & 使用者着急、旁观者催促 \\
        \bottomrule
    \end{tabular}
    \label{tab:noise}
\end{table}

\section{数据规模}

默认配置包含 20 个空间场景、3 个描述层级和 5 类干扰条件，因此单模型单轮采样形成：

\begin{equation}
20 \times 3 \times 5 = 300
\label{eq:prompt-size}
\end{equation}

条 Prompt。若比较 3 个模型并重复采样 3 次，则总回复量可达到 $300 \times 3 \times 3 = 2700$ 条。该规模能够支持准确率、翻转率、聚类和关联规则等多种数据挖掘分析。
